\documentclass[../../main.tex]{subfiles}
\begin{document}
\section{向量空间的定向}\label{section:ch15-vector-spaces}

\begin{figure}[htbp]
\centering
\begin{minipage}{0.48\textwidth}
\centering
\includegraphics[width=\linewidth]{fig15-1-sphere.png}
\caption{球面是可定向的}
\label{figure:ch15-sphere}
\end{minipage}
\end{figure}

\begin{figure}[htbp]
\centering
\begin{minipage}{0.48\textwidth}
\centering
\includegraphics[width=\linewidth]{fig15-2-mobius.png}
\caption{Möbius 带是不可定向的}
\label{figure:ch15-mobius}
\end{minipage}
\end{figure}

我们从向量空间的定向开始。我们都熟悉某些挑选 $\mathbb{R}^1,\mathbb{R}^2,\mathbb{R}^3$ 优选有序基的非正式规则（见图\ref{figure:ch15-bases}）。我们通常为 $\mathbb{R}^1$ 选一个指向右方（即沿正方向）的基。$\mathbb{R}^2$ 的一族自然的优选有序基，是那些从第一个向量到第二个向量的旋转为逆时针方向的基。每个学过向量微积分的学生都遇到过 $\mathbb{R}^3$ 中的“右手”基：它们是这样的有序基 $(E_1,E_2,E_3)$，当你右手的手指从 $E_1$ 弯曲到 $E_2$ 时，拇指指向 $E_3$ 的方向。

虽然“朝右”、“逆时针”、“右手”并不是数学术语，但很容易把挑选 $\mathbb{R}^1,\mathbb{R}^2,\mathbb{R}^3$ 优选基的规则翻译成严格的数学语言：你可以验证，在这三种情形下，优选基都是那些从标准基过渡过去的过渡矩阵行列式为正的基。

在一个没有典范基的抽象向量空间中，我们再也无法判断哪些基是“正确定向的”。例如，若 $V$ 是单实变量、次数至多为 $2$ 的多项式所成的向量空间，谁又能说有序基 $(1,x,x^2)$ 与 $(x^2,x,1)$ 中哪个是“右手的”？一般而言，我们所能说的只是两个基具有“相同定向”意味着什么。

于是我们被引导引入下面的定义。设 $V$ 是一个维数 $n\geqslant1$ 的实向量空间。若 $V$ 的两个有序基 $(E_1,\dots,E_n)$ 与 $(\widetilde E_1,\dots,\widetilde E_n)$ 的过渡矩阵 $(B^j_i)$——由
\begin{align}
E_i=B^j_i\widetilde E_j,\label{align:ch15-transition}
\end{align}
定义——行列式为正，我们就说它们是\emph{一致定向的}。

\begin{exercise}\label{exercise:ch15-01}
证明：一致定向是 $V$ 的所有有序基所成集合上的一个等价关系，并证明恰有两个等价类。
\end{exercise}

若 $\dim V=n\geqslant1$，我们把 $V$ 的一个定向定义为有序基的一个等价类。若 $(E_1,\dots,E_n)$ 是 $V$ 的任意一个有序基，我们把它所确定的定向记作 $[E_1,\dots,E_n]$，而把相反定向记作 $-[E_1,\dots,E_n]$。一个向量空间连同其定向的选取，称为一个\emph{有向向量空间}。若 $V$ 已定向，则任何属于给定定向的有序基 $(E_1,\dots,E_n)$ 称为定向的或正定向的；任何不属于给定定向的基称为负定向的。

对于零维向量空间 $V$ 的特殊情形，我们把 $V$ 的一个定向简单定义为在两个数 $\pm1$ 中选其一。

\begin{example}\label{example:ch15-standard-orientation}
$\mathbb{R}^n$ 由标准基所确定的定向 $[e_1,\dots,e_n]$ 称为\emph{标准定向}。你应使自己确信：在我们通常用图形表示坐标轴的方式下，$\mathbb{R}^1$ 的一个定向基就是指向右方的基；$\mathbb{R}^2$ 的一个定向基就是从第一个基向量到第二个基向量的旋转为逆时针的基；$\mathbb{R}^3$ 的一个定向基就是右手基。（这些可作为“朝右”、“逆时针”、“右手”等词的数学定义。）$\mathbb{R}^0$ 的标准定向定义为 $+1$。
\end{example}

定向与斜对称张量之间有一个重要联系，表述于下一个命题中。

\begin{proposition}\label{proposition:ch15-n-covector-orientation}
设 $V$ 是一个 $n$ 维向量空间。每个非零元素 $\omega\in\Lambda^n(V^*)$ 都按如下方式确定 $V$ 的一个定向 $\mathcal O_\omega$：若 $n\geqslant1$，则 $\mathcal O_\omega$ 是那些使 $\omega(E_1,\dots,E_n)>0$ 的有序基 $(E_1,\dots,E_n)$ 所成的集合；而若 $n=0$，则当 $\omega>0$ 时 $\mathcal O_\omega=+1$，当 $\omega<0$ 时 $\mathcal O_\omega=-1$。两个非零 $n$-余向量确定相同定向，当且仅当其中一个是另一个的正倍数。
\end{proposition}

\begin{proof}
零维情形是显然的，因为 $\Lambda^0(V^*)$ 的一个非零元素只不过是一个非零实数。因此可设 $n\geqslant1$。令 $\omega$ 是 $\Lambda^n(V^*)$ 的一个非零元素，$\mathcal O_\omega$ 表示那些 $\omega$ 取正值的有序基所成的集合。我们需证 $\mathcal O_\omega$ 恰为一个等价类。

设 $(E_i)$ 与 $(\widetilde E_j)$ 是 $V$ 的任意两个有序基，令 $B:V\to V$ 为把 $E_j$ 映到 $\widetilde E_j$ 的线性映射。这意味着 $\widetilde E_j=BE_j=B^i_jE_i$，故 $B$ 就是这两个基之间的过渡矩阵。由命题\ref{proposition:ch14-n-covector-det}，
\begin{align*}
\omega(\widetilde E_1,\dots,\widetilde E_n)
=\omega(BE_1,\dots,BE_n)=(\det B)\,\omega(E_1,\dots,E_n).
\end{align*}
由此可知，基 $(\widetilde E_j)$ 与 $(E_i)$ 一致定向，当且仅当 $\omega(E_1,\dots,E_n)$ 与 $\omega(\widetilde E_1,\dots,\widetilde E_n)$ 同号，这等于说 $\mathcal O_\omega$ 是一个等价类。最后一个结论随后容易得到。
\end{proof}

\begin{figure}[htbp]
\centering
\includegraphics[width=0.7\textwidth]{fig15-3-bases.png}
\caption{$\mathbb{R}^1,\mathbb{R}^2,\mathbb{R}^3$ 的定向基}
\label{figure:ch15-bases}
\end{figure}

若 $V$ 是一个已定向的 $n$ 维向量空间，$\omega$ 是一个按上述命题方式确定 $V$ 定向的 $n$-余向量，我们就说 $\omega$ 是一个（正）定向的 $n$-余向量。例如，$n$-余向量 $e^1\wedge\cdots\wedge e^n$ 对于 $\mathbb{R}^n$ 的标准定向是正定向的。

对任意 $n$ 维向量空间 $V$，空间 $\Lambda^n(V^*)$ 是一维的。命题\ref{proposition:ch15-n-covector-orientation}表明，为 $V$ 选取一个定向，等价于在 $\Lambda^n(V^*)\setminus\{0\}$ 的两个连通分支中选其一。这一表述对零维向量空间也适用，并解释了我们为何以那种方式定义零维空间的定向。
\end{document}
